\documentclass[11pt]{article}
\usepackage[margin=1in]{geometry}
\usepackage{microtype,booktabs,longtable,array,amsmath,amssymb,url,graphicx}
\usepackage[hidelinks]{hyperref}
\newcommand{\sgaeia}{\textsc{SGAEIA}}
\newcommand{\repo}{\url{https://github.com/aridiosilva/SGAEIA}}
\author{Aridio Silva\\\texttt{@aridiosilva}\\Independent Researcher\\\url{https://github.com/aridiosilva/SGAEIA}}
\date{September 2026}
\title{Governed Autonomy for Distributed Agentic Edge AI: Identity, Delegation, Zero Trust, Continuous GRC, and Formal Security Invariants}
\begin{document}\maketitle
\begin{abstract}
Distributed agentic AI creates authority-bearing non-human actors that can plan, invoke tools, delegate, use persistent memory, and operate under partial connectivity. This paper formalizes SGAEIA's governed-autonomy security model: workload identity, capability-based authorization, bounded delegation, criticality/autonomy classes, Agent Trust Zones, risk-adaptive authority, policy decision/enforcement separation, revocation, evidence-as-code, and fail-secure edge operation. The model is an executable reference specification, not a production certification.
\end{abstract}
\noindent\textbf{Project:} SGAEIA --- Secure Governed Autonomous Edge Intelligence Architecture.\\

\section{Motivation}
Edge intelligence places AI close to data and physical processes \cite{zhou2019edge}; agentic systems add tool use, memory, planning, and collaboration \cite{tran2025}. Security must therefore govern the authority-bearing agent, not only the model endpoint. Emerging agentic threat research highlights tool misuse, memory poisoning, autonomy, and coordination risks \cite{datta2025,owasp,mitreatlas}.

\section{Security Objective}
The objective is
\[
UsefulAutonomy\land BoundedAuthority\land ExplicitTrust\land Revocability\land Auditability.
\]
No model-generated semantic judgment alone is sufficient to authorize a critical action.

\section{Identity and Capability}
Every governed agent receives a non-human identity bound to its workload/execution context. SPIFFE demonstrates a standard model for verifiable workload identity \cite{spiffe}. Identity is necessary but insufficient: authorization also depends on capability, policy, risk, context, and delegation.

\section{Delegation}
For parent $p$ and child $c$,
\[
C(c)\subseteq C_{delegated}(p),\quad depth(c)\leq MaxDepth.
\]
Delegation also carries expiry, resource, purpose, and context constraints. Authority cannot be created merely by an agent claiming it.

\section{Criticality, Autonomy, and Trust Zones}
Criticality $L0\ldots L4$ and autonomy $A0\ldots A5$ are independent dimensions. The reference invariant is $L4\land A5\Rightarrow DENY$. Effective autonomy is the minimum permitted by configured, risk, context, trust, and policy bounds. ATZ-0 through ATZ-5 classify control requirements without granting transitive trust.

\section{Authorization and Zero Trust}
For agent $a$, request $r$, object $o$, context $c$, and time $t$:
\[
Authorize=I(a)\land Cap(a,o)\land P(a,r,o,c)\land R(a,r,o,c,t)<\theta\land T(a)>\tau.
\]
The PDP/PEP is external to the LLM. NIST Zero Trust supports identity-centric, explicit policy enforcement rather than network-location trust \cite{nist207,nist207a}.

\section{Fail-Secure Edge Operation}
Edge disconnection must not increase authority:
\[
OfflineMode\Rightarrow Authority_{offline}\leq Authority_{online}.
\]
Critical operations may require cached bounded policy, short-lived identity, human approval, independent safety interlocks, or denial.

\section{Revocation}
Revocation is temporal:
\[
RevocationRequested(t_0)\Rightarrow AuthorityUnavailable(t\leq t_0+\Delta_{max}).
\]
The bound is deployment-specific and must be measured across identity, policy, mesh, caches, tools, and edge nodes.

\section{Continuous GRC}
NIST AI RMF and its GenAI Profile motivate lifecycle governance \cite{nistrmf,nistgenai}. SGAEIA turns the governance chain into runtime evidence:
\[
Requirement\rightarrow Risk\rightarrow Control\rightarrow Policy\rightarrow Enforcement\rightarrow Evidence.
\]

\section{Threat-to-Control Mapping}
\begin{longtable}{p{.28\linewidth}p{.62\linewidth}}\toprule Threat&Representative controls\\\midrule
Agent impersonation&workload identity, attestation, short-lived credentials, revocation\\
Delegation escalation&subset/depth/expiry/context checks\\
Prompt/tool injection&untrusted-input zone, tool mediation, capability checks\\
Memory/RAG poisoning&provenance, source policy, isolation, retrieval authorization\\
Autonomous lateral movement&mesh identity, egress control, least privilege\\
Unsafe actuation&criticality limits, approval, independent safety barrier\\
Evidence suppression&buffering/sealing/correlation/monitoring\\\bottomrule\end{longtable}

\section{Security Invariants}
\begin{align*}
AgentWithoutIdentity&\Rightarrow DENY\\
RevokedAgent&\Rightarrow DENY\\
DelegationDepth>MaxDepth&\Rightarrow DENY\\
ChildCapabilities\nsubseteq ParentDelegatedCapabilities&\Rightarrow DENY\\
AttestationFailed\land CriticalAgent&\Rightarrow STOP\lor QUARANTINE.
\end{align*}

\section{Formalization and Limits}
TLA+ and Alloy artifacts are included as specifications; selected equivalent finite-state properties are tested in Python. No claim is made that TLA+/Alloy model checking has been executed. Risk scoring, safety constraints, and revocation bounds require domain-specific and live validation.

\section{Conclusion}
Agentic Edge AI requires security to govern authority rather than connectivity alone. SGAEIA makes identity, capability, delegation, risk, enforcement, evidence, and revocation explicit and testable.
\section*{Availability}
Open reference implementation and specifications: \repo.
\begin{thebibliography}{99}
\bibitem{shi2016edge} W. Shi et al., ``Edge Computing: Vision and Challenges,'' IEEE Internet of Things Journal, 3(5), 637--646, 2016. doi:10.1109/JIOT.2016.2579198.
\bibitem{zhou2019edge} Z. Zhou et al., ``Edge Intelligence: Paving the Last Mile of Artificial Intelligence With Edge Computing,'' Proceedings of the IEEE, 107(8), 1738--1762, 2019. doi:10.1109/JPROC.2019.2918951.
\bibitem{li2018} E. Li, Z. Zhou, X. Chen, ``Edge Intelligence: On-Demand Deep Learning Model Co-Inference with Device-Edge Synergy,'' MECOMM, 2018. doi:10.1145/3229556.3229562.
\bibitem{nist207} NIST, Zero Trust Architecture, SP 800-207, 2020. doi:10.6028/NIST.SP.800-207.
\bibitem{nist207a} NIST, A Zero Trust Architecture Model for Access Control in Cloud-Native Applications in Multi-Cloud Environments, SP 800-207A, 2023. doi:10.6028/NIST.SP.800-207A.
\bibitem{nistrmf} NIST, Artificial Intelligence Risk Management Framework (AI RMF 1.0), AI 100-1, 2023. doi:10.6028/NIST.AI.100-1.
\bibitem{nistgenai} C. Autio et al., Artificial Intelligence Risk Management Framework: Generative Artificial Intelligence Profile, NIST AI 600-1, 2024. doi:10.6028/NIST.AI.600-1.
\bibitem{nistssdf} NIST, Secure Software Development Practices for Generative AI and Dual-Use Foundation Models, SP 800-218A, 2024.
\bibitem{mitreatlas} MITRE, ATLAS: Adversarial Threat Landscape for Artificial-Intelligence Systems, \url{https://atlas.mitre.org/}, accessed Sep. 2026.
\bibitem{owasp} OWASP GenAI Security Project, Agentic Security Initiative and Agent Control Standard, \url{https://genai.owasp.org/}, accessed Sep. 2026.
\bibitem{spiffe} SPIFFE Project, The SPIFFE Standard, \url{https://spiffe.io/docs/latest/spiffe-specs/}, accessed Sep. 2026.
\bibitem{etsi} ETSI, Multi-access Edge Computing (MEC), \url{https://www.etsi.org/technical-groups/mec/}, accessed Sep. 2026.
\bibitem{openapi} OpenAPI Initiative, OpenAPI Specification 3.2.0, 2025.
\bibitem{asyncapi} AsyncAPI Initiative, AsyncAPI Specification 3.1.0, 2026.
\bibitem{tran2025} K.-T. Tran et al., ``Multi-Agent Collaboration Mechanisms: A Survey of LLMs,'' arXiv:2501.06322, 2025.
\bibitem{datta2025} S. Datta et al., ``Agentic AI Security: Threats, Defenses, Evaluation, and Open Challenges,'' arXiv:2510.23883, 2025.
\end{thebibliography}
\end{document}
